\documentclass[11pt]{article}

% ============================================================
% Milestone 1 -- Project Proposal & Literature Review Template
% DIP Team Project
% ============================================================
\usepackage[margin=1in]{geometry}
\usepackage[T1]{fontenc}
\usepackage[utf8]{inputenc}
\usepackage{amsmath,amssymb}
\usepackage{graphicx}
\usepackage{booktabs}
\usepackage[round]{natbib}
\usepackage{hyperref}
\hypersetup{colorlinks=true, linkcolor=blue!50!black, citecolor=blue!50!black, urlcolor=blue!50!black}
\usepackage{enumitem}
\usepackage{xcolor}

\newcommand{\todo}[1]{\textcolor{red}{\textbf{[TODO: #1]}}}

\title{\LARGE Project Proposal \& Literature Review \\[4pt]
\large \todo{Project Title}}
\author{
  \todo{Student A Full Name} \and
  \todo{Student B Full Name} \and
  \todo{Student C Full Name}
}
\date{Milestone 1 --- \todo{Date} \\{} [COURSE CODE] --- Digital Image Processing}

\begin{document}
\maketitle

\begin{abstract}
\todo{3--5 sentences: what problem are you tackling, why does it matter, what is your
planned approach, and what public dataset/baseline will you use. Write this last.}
\end{abstract}

\section{Problem Statement \& Motivation}
\label{sec:motivation}
\todo{What image-processing problem are you solving? Why is it interesting or useful?
How does it connect to material from the course (name the week/topic)? 1 paragraph.}

\section{Related Work}
\label{sec:related}
\todo{Summarize at least 4 references, at least 2 of which should be from 2023--2026.
Use the table below as a starting organizer, then write 1--2 paragraphs of prose
connecting them to your plan.}

\begin{table}[h]
\centering
\caption{Key references.}
\label{tab:refs}
\begin{tabular}{@{}p{0.30\linewidth}p{0.12\linewidth}p{0.45\linewidth}@{}}
\toprule
\textbf{Reference} & \textbf{Year} & \textbf{Relevance to this project} \\
\midrule
\citep{example2024} & 2024 & \todo{one sentence} \\
\citep{example2018}  & 2018 & \todo{one sentence} \\
\bottomrule
\end{tabular}
\end{table}

\section{Dataset \& Baseline Plan}
\label{sec:data}
\todo{Name the public dataset(s), give exact size/subset you will use, and the license.
State exactly which baseline (classical or prior learned method) you will reproduce
and compare against.}

\begin{table}[h]
\centering
\caption{Dataset summary.}
\label{tab:data}
\begin{tabular}{@{}lll@{}}
\toprule
\textbf{Dataset} & \textbf{Subset used} & \textbf{Link} \\
\midrule
\todo{name} & \todo{e.g., 800 train / 100 val crops} & \todo{URL} \\
\bottomrule
\end{tabular}
\end{table}

\section{Proposed Method}
\label{sec:method}
\todo{Sketch your planned approach: architecture diagram or pseudocode, what is novel
or extended relative to the baseline, and why you expect it to work.}

\section{Feasibility \& Compute Plan}
\label{sec:feasibility}
\todo{State your compute tier (CPU-OK / GPU-light), expected training time per
experiment, and confirm this fits a laptop-class machine within the project window.}

\subsection{Minimum Viable Deliverable (MVP)}
\todo{The smallest version of this project that still produces a valid, reportable
comparison against the baseline.}

\subsection{Stretch Goal}
\todo{What you will attempt if the MVP is on track by Week 5--6.}

\section{Risk Assessment}
\label{sec:risk}
\begin{itemize}[leftmargin=*]
  \item \todo{Risk 1 and its mitigation / fallback.}
  \item \todo{Risk 2 and its mitigation / fallback.}
\end{itemize}

\section{Work Breakdown}
\label{sec:workbreakdown}
\begin{table}[h]
\centering
\caption{Milestone task allocation (see also the project catalog entry for this
project).}
\label{tab:tasks}
\begin{tabular}{@{}p{0.12\linewidth}p{0.26\linewidth}p{0.26\linewidth}p{0.26\linewidth}@{}}
\toprule
\textbf{Milestone} & \textbf{Student A} & \textbf{Student B} & \textbf{Student C} \\
\midrule
M1 (Wk 3) & \todo{task} & \todo{task} & \todo{task} \\
M2 (Wk 6) & \todo{task} & \todo{task} & \todo{task} \\
M3 (Wk 9) & \todo{task} & \todo{task} & \todo{task} \\
\bottomrule
\end{tabular}
\end{table}

\bibliographystyle{unsrtnat}
\bibliography{references}

\end{document}
